\documentclass[11pt]{article}
\usepackage{amsmath,amssymb}
\usepackage[hidelinks]{hyperref}
\title{When a stopped factory station retains a job}
\date{}
\begin{document}
\maketitle

The official RoboCup Logistics League archive named \texttt{rcll100.gz}
contains 100 reports and 900 orders, but no realized fault history:
there is one initial idle entry per machine, no recorded positive-duration
downtime and no completed deliveries. It supplies nominal benchmark
inputs, not independent evidence of fault recovery. A separate simulation
based on the league rules tests an actual recovery distinction. A paused
station can resume its old work; after confirmed improper use the
referee removes affected material. A controller that sees only that
a station is unavailable cannot tell which product should enter
the last early slot at an alternate station.

\section{What the published archive supports}
The \href{https://github.com/robocup-logistics/rcll-refbox/tree/master/benchmarks}{official
RefBox repository} publishes the archive. This audit pins its Git blob
\texttt{b54f8c7ba3f6ee4765d49f99013466275009d414}
and its compressed SHA-256 in the replay. The eligibility test was set
before looking at any report records. An actual event-level compound
case must establish a physically live material link, fault cause,
credible event ordering, the actor's information before the action,
and relevant workload and deadlines. Nominal geography or order
parameters alone do not meet that test.

A bounded, read-only parser walked the original
\href{https://pkg.go.dev/github.com/mongodb/mongo-tools/common/archive}{MongoDB
archive format}; it did not execute or restore the dump. It recovered
one collection, \texttt{rcll.game\_report}, with 100 reports. All 1,400
machine-history entries occur at game time zero and are idle.
Each report has exactly one entry per named machine. None of the
1,400 \texttt{down-period} pairs describes a nonnegative start and
later end, none of the broken reasons is populated, and no recorded
order unit was delivered. The 99 end-time metadata fields do not
supply the missing executions. Several key names look like trace
fields, but checking their values leaves zero eligible cases. This
is a finding about one pinned archive, not a claim that the league
does not keep detailed game reports or that no other public export
would work. It does not spend the enclosed-space mission holdout or
certify transfer of a frozen field controller.

\section{Different physical consequences of unavailability}
The \href{https://github.com/robocup-logistics/rcll-rulebook/blob/master/rulebook.tex}{official
rulebook}, pinned in the README, states that scheduled cap or ring
station downtime can last 30--60 seconds and that a product already
processing resumes after that delay. Confirmed improper use can
lead to the broken state and removal of material from processing,
input and output; the station remains unavailable for 30 seconds.
The same material-removal inference does not automatically follow
from a generic hardware anomaly or any reset. The rulebook
distinguishes technical recovery and misuse, and its machine lights
distinguish scheduled down from broken. Broken state alone does
not certify that improper use, rather than a manual reset, was the
cause. A signal used below therefore either comes from direct
validation of the retained item or has an independent cause
certificate within the declared two-case envelope. The replay
does not assert working telemetry from the published reports.

In the separate development fixture, station X contains the unfinished
item \(P\) at the first fault. The requirement is for durably
collected \(P\) and \(Q\) by a fixed horizon, not a referee score.
Under scheduled pause, X retains \(P\) and resumes.
Under referee removal it loses that item, the order for \(P\)
still stands, and X can instead start a prepared \(Q\) on reopening.
Station Y can start one prequalified order now, but a second
interruption arrives before the plant can carry another order.
All launch and collection authority in this model is declared.
A spare base for the possible rebuild is staged before the test;
the model does not give the controller a free clone of the resident.
The first choice uses the same capacities, time and materials in either
world. Processing and handling durations are synthetic, not archive
measurements.

Both items arrive if Y starts \(Q\) in the pause case or \(P\) in
the removal case. The other first order yields attempts at
the same letter while the other required item misses the deadline.
A policy that sees only unavailability cannot guarantee both.
A timely, source-validated retention bit enables both.
There is no fall-back that bypasses the missing information within
the specified one-launch action class. A direct planner with the
same bit and capacities also succeeds.

\section{Completion versus secure collection}
The action consumes work time \(p\) and a declared handling interval
\(\ell\). The machine's output can be ready before the item is collected
outside the interruption's physical scope. If launch is \(s\), the
uninterrupted secure time would be \(c=s+p+\ell\). Let the second
fault occur at \(f\), the deadline be \(T\), and \(\delta=0\) when
a simultaneous pickup precedes the fault and 1 when the fault is
first. The pre-fault arm succeeds exactly when
\[
       c\leq\min(f-\delta,T).
\]
If the second event is a pause of duration \(D\), unfinished
processing or handling resumes with a delay \(D\), under the
specified same-station handling law. The full success test is the
union of the pre-fault arm and \(c\leq T-D\). For destructive removal,
only the pre-fault arm remains; already completed but uncollected
material can be removed. The law for scheduled delay during
handling is a fixture assumption, not an additional official rule.

In the replay, a product is ready at tick 28, but a fault-first event
at tick 30 destroys the uncollected product. Collection-first at
the same timestamp preserves it. Another test delays observation
one tick and uses a scheduled second pause; the selected transfer
then collects at tick 60 and fails a deadline of 59. Four Lean
statements check the arithmetic, equal-time distinction and
opposite-first-choice obstruction. The event model separately
tracks a running job, output-ready material, paused occupancy,
physical removal and secure collection. All 82,944 cases agree
with an independent interval calculation.

\section{Evidence boundary}
\href{https://doi.org/10.1016/S0377-2217(00)00322-2}{Lee and Lin}
already distinguish resuming and nonresuming jobs in machine
scheduling. The additional interface here is actual retained
material and the witness that crosses the boundary. It is not
a new general scheduling method. The event model and the archive
audit answer different questions. The latter does not confirm
the former. A future independent case needs matched material identity,
occupancy and collection events, order and fault evidence,
timestamps, what each agent saw before committing, and an unchanged
comparison under the same physical resources. An independently
chosen layout or order list is not yet that case. This is I1
development evidence, not a decision on inventory-wide applicability
or a claim that a stronger established planner loses.
\end{document}
